\authorchapter{Einleitung}{Ereny Micail}

Wer kennt das nicht: In der Schule oder im Team läuft alles über verschiedene Apps.
Aufgaben hier, Nachrichten dort, Dateien woanders.
Nach ein paar Wochen weiß keiner mehr, was Sache ist.
Das nennt man \textit{Context Switching} -- und es kostet nicht nur Zeit, sondern
sorgt auch dafür, dass Infos verloren gehen und keiner mehr weiß, wer für was
zuständig ist.

Genau das wollten wir mit \textit{EffSync} lösen.
\textit{EffSync} ist eine Webanwendung, die Aufgabenverwaltung und Chat in einem
System zusammenbringt -- alles an einem Ort, kein Hin- und Herwechseln mehr.
Nutzer können Tasks anlegen, zuweisen und verfolgen, und gleichzeitig direkt im
Projekt-Kontext miteinander schreiben.
Das Ziel der Arbeit war, genau das umzusetzen und zu zeigen, dass es funktioniert.
Dabei haben wir drei Dinge in den Mittelpunkt gestellt: eine saubere Task-Struktur
mit klaren Berechtigungen, Echtzeit-Kommunikation über WebSocket und eine ordentliche
Qualitätssicherung durch Tests und einen Usability-Test.

Wir haben uns auf bekannte Konzepte gestützt – etwa die Modularisierung nach
Parnas~\cite{parnas1972criteria}, die Grundlagen der Softwarearchitektur nach
Bass et al.\ \cite{bass2012software} sowie Scrum nach Schwaber und
Sutherland~\cite{schwaber2020scrum} – und zusätzlich auf die Dokumentationen der verwendeten
Technologien: REST, WebSocket, Prisma und PostgreSQL.
Auch der Datenschutz spielte eine Rolle; dabei orientierten wir uns an der DSGVO~\cite{dsgvo2016}
und den Empfehlungen der ENISA~\cite{enisa2019pseudonymisation}.
Für die Projektdurchführung setzten wir auf Scrum: Vue.js im Frontend sowie NestJS mit
Prisma und PostgreSQL im Backend.
Zunächst entwickelten wir das Datenmodell und die zentrale Anwendungslogik, anschließend
ergänzten wir Echtzeit-Funktionalitäten und Optimierungen der Bedienbarkeit.
Die Qualität stellten wir durch Unit-Tests, End-to-End-Tests und einen moderierten
Usability-Test sicher.
Was diese Arbeit nicht macht: native Mobile-Apps, Mandantenfähigkeit für mehrere
Organisationen, Hochskalierung für sehr viele Nutzer oder vollständige Offline-Nutzung.
Das wäre alles wichtig für ein echtes Produkt, sprengt aber den Rahmen einer
einsemestrigen Diplomarbeit.
Diese Punkte haben wir in Kapitel~\ref{ch:ausblicke} als Ideen für die Zukunft notiert.

So ist die Arbeit aufgebaut:
Kapitel~\ref{ch:ausgangslage} erklärt, wo wir gestartet sind und was das Ziel war.
Kapitel~\ref{ch:projektmanagement} und~\ref{ch:workflow} zeigen, wie wir das Projekt
organisiert und mit Git gearbeitet haben.
Kapitel~\ref{ch:projektbeschreibung} stellt \textit{EffSync} und die Anforderungen
vor, bevor Kapitel~\ref{ch:systemarchitektur}
und~\ref{ch:implementierung_kernfunktionen} erklären, wie das System aufgebaut ist
und wie wir die wichtigsten Funktionen umgesetzt haben.
Kapitel~\ref{ch:testing_quality} beschreibt, wie wir die Qualität gesichert haben.
Kapitel~\ref{ch:ergebnisse} und~\ref{ch:ausblicke} zeigen, was rausgekommen ist
und was man noch machen könnte.
Am Ende steht in Kapitel~\ref{ch:reflexion} unsere Reflexion -- eine gemeinsame
und zwei persönliche.